%
% listener_lifecycle_nowait_buggy.tex
%
% FIDELITY CONTROL for docs/listener_lifecycle.tex: the ROUND-2 design (retire
% on timeout, one untyped reply sink per generation, one request in flight
% per link) in the presence of the replicator's un-awaited send_scene.  It is
% the round-2 fixed specification verbatim plus ONE operation, SendNoWait: a
% write that makes a reply owed and has no waiter.  Round-2's retire rule
% assumed every write that elicits a reply is awaited under the request lock;
% the display acks every SceneMessage, so that assumption is false.
%
% Expected: L1 to L5 hold; L6 VIOLATED
%   (Dial; SendNoWait; Send; Await; StartListener(ta); CheckPermit(ta);
%    Feed(1,ta); RouteLive(1,ta); Consume(1))
% -- no timeout is involved: a later show() takes the ack of an earlier
% fire-and-forget scene.
%
% Goals (negations of the invariants; see listener_lifecycle.tex, round 2):
%   L1 (qrep /= {} & (held = fclear or {r | r:qrep & rgen(r) /= curGen} /= {}))
%        or (cstate = cwait & (held = fclear or rgen(creq) /= curGen))
%   L2 held = fclear & curGen > 0 & curGen /: (broken \/ deliberate)
%   L3 afterStop(ta) > 1 or afterStop(tb) > 1
%   L4a lst(ta) /= lexit & lst(tb) /= lexit & lgen(ta) = lgen(tb)
%   L4b {r | r:bufr(ta) & rgen(r) /= lgen(ta)} /= {}
%        or {r | r:bufr(tb) & rgen(r) /= lgen(tb)} /= {}
%   L5 crossGen = fset      L6 crossReq = fset
% Check:  fuzz FILE ; probcli FILE -model_check -nodead -p DEFAULT_SETSIZE 2 \
%           -p MAX_OPERATIONS 200 -goal "<negation>"
%
\documentclass[a4paper,10pt,fleqn]{article}
\usepackage[margin=1in]{geometry}
\usepackage{fuzz}
\usepackage[colorlinks=true,linkcolor=blue,citecolor=blue,urlcolor=blue]{hyperref}

\begin{document}

\title{Listener Lifecycle: the Round-2 Design with an Unawaited Send}
\author{Fidelity control for listener\_lifecycle.tex}
\date{October 2026}
\maketitle

% ============================================================
\section{Constants and Free Types}
% ============================================================

Three generations: the original, the replacement, and one more, so that a
replacement can itself be replaced while an abandoned reader of the first is
still alive.  Two requests: the least that lets a late reply meet a later
request.  One buffered reply is the least that lets an abandoned reader have
something left to route.

\begin{axdef}
  maxGen : \nat \\
  maxReq : \nat \\
  maxBuf : \nat
\where
  maxGen = 3 \\
  maxReq = 2 \\
  maxBuf = 1
\end{axdef}

\begin{zed}
  T ::= ta | tb
\end{zed}

\begin{zed}
  LST ::= lexit | lrun | lstop
\end{zed}

\begin{zed}
  FSTATE ::= fclear | fset
\end{zed}

\begin{zed}
  CST ::= cidle | csent | cwait
\end{zed}

\texttt{T}: the two reader threads, as slots.  \texttt{LST}: a slot is idle
(\texttt{lexit}), running with its stopping flag clear (\texttt{lrun}), or
alive with the flag set (\texttt{lstop} --- \texttt{stop} timed out, or has
not yet joined: the abandoned thread).  \texttt{CST}: the caller's request is
idle, written and not yet awaited (\texttt{csent}), or awaited
(\texttt{cwait}).

% ============================================================
\section{State}
% ============================================================

The state is split into four groups, then joined, so that an operation can
say what it leaves alone.

\begin{schema}{Conn}
  curGen : \nat \\
  held : FSTATE
\where
  curGen \leq maxGen \\
  held = fset \implies curGen > 0
\end{schema}

\begin{schema}{Req}
  nreq : \nat \\
  creq : \nat \\
  cstate : CST \\
  rgen : (0 \upto maxReq) \fun (0 \upto maxGen) \\
  owed : \power \nat \\
  qrep : \power \nat
\where
  nreq \leq maxReq \\
  creq \leq maxReq \\
  owed \subseteq 1 \upto maxReq \\
  qrep \subseteq 1 \upto maxReq
\end{schema}

\begin{schema}{Thr}
  lst : T \fun LST \\
  lgen : T \fun \nat \\
  armed : T \fun FSTATE \\
  bufr : T \fun (\power \nat) \\
  afterStop : T \fun \nat
\where
  \forall t : T \spot lgen~t \leq maxGen \land \# (bufr~t) \leq maxBuf
    \land bufr~t \subseteq 1 \upto maxReq \land afterStop~t \leq maxBuf + 1
\end{schema}

\begin{schema}{Hist}
  broken : \power \nat \\
  deliberate : \power \nat \\
  crossGen : FSTATE \\
  crossReq : FSTATE
\where
  broken \subseteq 0 \upto maxGen \\
  deliberate \subseteq 0 \upto maxGen
\end{schema}

\begin{schema}{LL}
  Conn \\
  Req \\
  Thr \\
  Hist
\end{schema}


\texttt{curGen} is the latest generation dialed (zero before the first dial);
\texttt{held} says whether the connection holds that generation's socket ---
\texttt{is\_open}.  Every earlier generation is closed.

In \texttt{Req}, requests are numbered \texttt{1..nreq}; \texttt{rgen} gives
the generation each was written on; \texttt{creq} and \texttt{cstate} are the
caller's request and where it stands.  \texttt{owed} are the requests the
display may still answer; \texttt{qrep} are the replies sitting in sinks (a
reply is the number of its request).

In \texttt{Thr}, \texttt{lgen} is the generation a reader was started on;
\texttt{bufr} are the replies it has read and not yet routed;
\texttt{armed} is the permission the reader obtained the last time it looked
at its stopping flag while the flag was clear (the loop-top test in
\texttt{\_run}, and the re-test this design adds before a route or a report).
A step that routes or reports consumes the permission.  The flag can be set
\emph{between} the look and the step, which is why a permission is state: it
is the window, and the model keeps it open.  \texttt{afterStop} counts the
routing and reporting steps a reader has taken since its flag was set.

\texttt{Hist} holds history variables, never read by an operation guard, so
that invariants are statements about what happened.  \texttt{broken}: the
generations whose own reader found their own stream broken.
\texttt{deliberate}: the generations the link, or a timeout, closed.
\texttt{crossGen}: a caller consumed a reply of another generation.
\texttt{crossReq}: a caller consumed a reply to a request other than its own.

\section{Initialization}

\begin{schema}{Init}
  LL'
\where
  curGen' = 0 \\
  held' = fclear \\
  nreq' = 0 \\
  creq' = 0 \\
  cstate' = cidle \\
  \forall i : 0 \upto maxReq \spot rgen'~i = 0 \\
  owed' = \emptyset \\
  qrep' = \emptyset \\
  lst' = \{ ta \mapsto lexit, tb \mapsto lexit \} \\
  lgen' = \{ ta \mapsto 0, tb \mapsto 0 \} \\
  armed' = \{ ta \mapsto fclear, tb \mapsto fclear \} \\
  bufr' = \{ ta \mapsto \emptyset, tb \mapsto \emptyset \} \\
  afterStop' = \{ ta \mapsto 0, tb \mapsto 0 \} \\
  broken' = \emptyset \\
  deliberate' = \emptyset \\
  crossGen' = fclear \\
  crossReq' = fclear
\end{schema}


Nothing dialed, no request, no reader, no reply, nothing recorded.

% ============================================================
\section{Operations on the Connection and the Caller}
% ============================================================


\subsection{Dial}

\begin{schema}{Dial}
  \Delta LL \\
  \Xi Req \\
  \Xi Thr \\
  \Xi Hist
\where
  held = fclear \\
  curGen < maxGen \\
  curGen' = curGen + 1 \\
  held' = fset
\end{schema}
\texttt{DisplayConnection.open}: a new generation, a new socket.  Only a
closed connection dials.  The new generation has no sink content; nothing of
an earlier generation is carried over.

\subsection{Close}

\begin{schema}{Close}
  \Delta LL \\
  \Xi Thr
\where
  held = fset \\
  held' = fclear \\
  curGen' = curGen \\
  qrep' = \emptyset \\
  owed' = \{ r : owed | rgen~r \neq curGen \} \\
  nreq' = nreq \\
  creq' = creq \\
  rgen' = rgen \\
  (cstate = cwait \implies cstate' = cidle) \\
  (cstate \neq cwait \implies cstate' = cstate) \\
  deliberate' = deliberate \cup \{ curGen \} \\
  broken' = broken \\
  crossGen' = crossGen \\
  crossReq' = crossReq
\end{schema}
The link closes the connection on purpose: \texttt{DisplayLink.close},
\texttt{ClientRegistry.drop}.  The generation's sink is discarded, a waiting
caller fails, and the replies the display still owed on that socket are gone
with it.  Readers are not touched; \texttt{StopJoined} and
\texttt{StopTimeout} are separate steps, in either order.  The link-level
event queue is not touched (Decision~7).

\subsection{Send}

\begin{schema}{Send}
  \Delta LL \\
  \Xi Conn \\
  \Xi Thr \\
  \Xi Hist
\where
  cstate = cidle \\
  held = fset \\
  nreq < maxReq \\
  nreq' = nreq + 1 \\
  creq' = nreq + 1 \\
  rgen' = rgen \oplus \{ nreq + 1 \mapsto curGen \} \\
  owed' = owed \cup \{ nreq + 1 \} \\
  qrep' = qrep \\
  cstate' = csent
\end{schema}
The caller, holding the request lock, writes a request on the current socket.
Its reply is owed by the display.  The write returns the generation
(\texttt{rgen}) --- Decision~4.

\subsection{Await}

\begin{schema}{Await}
  \Delta LL \\
  \Xi Conn \\
  \Xi Thr \\
  \Xi Hist
\where
  cstate = csent \\
  held = fset \\
  rgen~creq = curGen \\
  nreq' = nreq \\
  creq' = creq \\
  rgen' = rgen \\
  owed' = owed \\
  qrep' = qrep \\
  cstate' = cwait
\end{schema}

\begin{schema}{AwaitFail}
  \Delta LL \\
  \Xi Conn \\
  \Xi Thr \\
  \Xi Hist
\where
  cstate = csent \\
  \lnot (held = fset \land rgen~creq = curGen) \\
  nreq' = nreq \\
  creq' = creq \\
  rgen' = rgen \\
  owed' = owed \\
  qrep' = qrep \\
  cstate' = cidle
\end{schema}
The caller waits on the sink of the generation it wrote on, if that
generation is still the live one; otherwise the wait fails at once
(\texttt{AwaitFail}).  A redial between \texttt{Send} and \texttt{Await} is
the case \texttt{AwaitFail} exists for.

\subsection{Consume}

\begin{schema}{Consume}
  \Delta LL \\
  \Xi Conn \\
  \Xi Thr \\
  r? : 1 \upto maxReq
\where
  cstate = cwait \\
  r? \in qrep \\
  rgen~r? = rgen~creq \\
  qrep' = qrep \setminus \{ r? \} \\
  cstate' = cidle \\
  nreq' = nreq \\
  creq' = creq \\
  rgen' = rgen \\
  owed' = owed \\
  broken' = broken \\
  deliberate' = deliberate \\
  crossGen' = crossGen \\
  (r? \neq creq \implies crossReq' = fset) \\
  (r? = creq \implies crossReq' = crossReq)
\end{schema}
The waiting caller takes a reply from its generation's sink.  The sink is
untagged: any reply in it will do, as in the code.  \texttt{crossReq} records
a reply that was not the caller's own; the fixed design never sets it,
because a sink holds only replies to requests written on its generation and
the generation is retired the moment a request is abandoned.

\subsection{TimeoutRetire}

\begin{schema}{TimeoutRetire}
  \Delta LL \\
  \Xi Thr
\where
  cstate = cwait \\
  held = fset \\
  rgen~creq = curGen \\
  held' = fclear \\
  curGen' = curGen \\
  qrep' = \emptyset \\
  owed' = \{ r : owed | rgen~r \neq curGen \} \\
  nreq' = nreq \\
  creq' = creq \\
  rgen' = rgen \\
  cstate' = cidle \\
  deliberate' = deliberate \cup \{ curGen \} \\
  broken' = broken \\
  crossGen' = crossGen \\
  crossReq' = crossReq
\end{schema}
The caller's budget expires while it waits.  The generation it wrote on is
ended: closed, its sink discarded, its owed replies gone.  This is the
timeout rule of Decision~5.  Nothing is reconnected here; the next
\texttt{Get} redials.

% ============================================================
\section{Operations on a Reader}
% ============================================================

Each takes the thread \texttt{t?}.

\subsection{StartListener}

\begin{schema}{StartListener}
  \Delta LL \\
  \Xi Conn \\
  \Xi Req \\
  \Xi Hist \\
  t? : T
\where
  lst~t? = lexit \\
  held = fset \\
  \forall u : T \spot lst~u \neq lrun \\
  \forall u : T \spot lst~u = lexit \lor lgen~u \neq curGen \\
  lst' = lst \oplus \{ t? \mapsto lrun \} \\
  lgen' = lgen \oplus \{ t? \mapsto curGen \} \\
  armed' = armed \oplus \{ t? \mapsto fclear \} \\
  bufr' = bufr \oplus \{ t? \mapsto \emptyset \} \\
  afterStop' = afterStop \oplus \{ t? \mapsto 0 \}
\end{schema}

\texttt{ListenerThread.start}.  A no-op while a running reader exists (the third line is the code's \texttt{if not stopping: return}).  The fourth line is this design's addition: a successor is started only on a generation no live thread is still reading.  Abandoning a reader means retiring its generation first (Decision~3).

\subsection{CheckPermit}

\begin{schema}{CheckPermit}
  \Delta LL \\
  \Xi Conn \\
  \Xi Req \\
  \Xi Hist \\
  t? : T
\where
  lst~t? = lrun \\
  armed~t? = fclear \\
  lst' = lst \\
  lgen' = lgen \\
  bufr' = bufr \\
  afterStop' = afterStop \\
  armed' = armed \oplus \{ t? \mapsto fset \}
\end{schema}

The reader looks at its stopping flag, finds it clear, and obtains permission for its next routing or reporting step.  Only a running reader can: once the flag is set, no permission is issued.

\subsection{Feed}

\begin{schema}{Feed}
  \Delta LL \\
  \Xi Conn \\
  \Xi Hist \\
  t? : T \\
  r? : 1 \upto maxReq
\where
  lst~t? = lrun \\
  held = fset \\
  lgen~t? = curGen \\
  r? \in owed \\
  rgen~r? = curGen \\
  \# (bufr~t?) < maxBuf \\
  owed' = owed \setminus \{ r? \} \\
  bufr' = bufr \oplus \{ t? \mapsto bufr~t? \cup \{ r? \} \} \\
  nreq' = nreq \\
  creq' = creq \\
  rgen' = rgen \\
  qrep' = qrep \\
  cstate' = cstate \\
  lst' = lst \\
  lgen' = lgen \\
  armed' = armed \\
  afterStop' = afterStop
\end{schema}

The reader reads a reply off its own generation's socket, while that socket is the open current one.  A reader over a closed or replaced socket reads nothing new; what it already holds stays in \texttt{bufr}. This is Decision~8: the reader holds its socket and never reads through \texttt{require()}.

\subsection{StopJoined}

\begin{schema}{StopJoined}
  \Delta LL \\
  \Xi Conn \\
  \Xi Req \\
  \Xi Hist \\
  t? : T
\where
  lst~t? = lrun \\
  lst' = lst \oplus \{ t? \mapsto lexit \} \\
  lgen' = lgen \\
  armed' = armed \oplus \{ t? \mapsto fclear \} \\
  bufr' = bufr \\
  afterStop' = afterStop
\end{schema}

\texttt{stop}: the flag is set and the join succeeds.

\subsection{StopTimeout}

\begin{schema}{StopTimeout}
  \Delta LL \\
  \Xi Conn \\
  \Xi Req \\
  \Xi Hist \\
  t? : T
\where
  lst~t? = lrun \\
  lst' = lst \oplus \{ t? \mapsto lstop \} \\
  lgen' = lgen \\
  armed' = armed \\
  bufr' = bufr \\
  afterStop' = afterStop \oplus \{ t? \mapsto 0 \}
\end{schema}

\texttt{stop}: the flag is set and the join times out.  The thread is alive, its permission (if any) intact, and it is abandoned.

\subsection{ExitStopped}

\begin{schema}{ExitStopped}
  \Delta LL \\
  \Xi Conn \\
  \Xi Req \\
  \Xi Hist \\
  t? : T
\where
  lst~t? = lstop \\
  lst' = lst \oplus \{ t? \mapsto lexit \} \\
  lgen' = lgen \\
  armed' = armed \oplus \{ t? \mapsto fclear \} \\
  bufr' = bufr \\
  afterStop' = afterStop
\end{schema}

The abandoned thread sees its flag at its next test and exits.

\subsection{RouteLive}

\begin{schema}{RouteLive}
  \Delta LL \\
  \Xi Conn \\
  \Xi Hist \\
  t? : T \\
  r? : 1 \upto maxReq
\where
  armed~t? = fset \\
  r? \in bufr~t? \\
  held = fset \\
  lgen~t? = curGen \\
  qrep' = qrep \cup \{ r? \} \\
  armed' = armed \oplus \{ t? \mapsto fclear \} \\
  bufr' = bufr \oplus \{ t? \mapsto bufr~t? \setminus \{ r? \} \} \\
  lst' = lst \\
  lgen' = lgen \\
  nreq' = nreq \\
  creq' = creq \\
  rgen' = rgen \\
  owed' = owed \\
  cstate' = cstate \\
  (lst~t? = lstop \implies
    afterStop' = afterStop \oplus \{ t? \mapsto afterStop~t? + 1 \}) \\
  (lst~t? \neq lstop \implies afterStop' = afterStop)
\end{schema}

A reader with permission routes one buffered reply into the sink of its own generation, when that generation is the live current one.

\subsection{RouteDead}

\begin{schema}{RouteDead}
  \Delta LL \\
  \Xi Conn \\
  \Xi Hist \\
  t? : T \\
  r? : 1 \upto maxReq
\where
  armed~t? = fset \\
  r? \in bufr~t? \\
  \lnot (held = fset \land lgen~t? = curGen) \\
  qrep' = qrep \\
  armed' = armed \oplus \{ t? \mapsto fclear \} \\
  bufr' = bufr \oplus \{ t? \mapsto bufr~t? \setminus \{ r? \} \} \\
  lst' = lst \\
  lgen' = lgen \\
  nreq' = nreq \\
  creq' = creq \\
  rgen' = rgen \\
  owed' = owed \\
  cstate' = cstate \\
  (lst~t? = lstop \implies
    afterStop' = afterStop \oplus \{ t? \mapsto afterStop~t? + 1 \}) \\
  (lst~t? \neq lstop \implies afterStop' = afterStop)
\end{schema}

If the generation has been closed or replaced, its sink is gone and the reply is dropped.  Nothing here reaches the current generation.

\subsection{BrokenOwn}

\begin{schema}{BrokenOwn}
  \Delta LL \\
  t? : T
\where
  armed~t? = fset \\
  held = fset \\
  lgen~t? = curGen \\
  held' = fclear \\
  curGen' = curGen \\
  qrep' = \emptyset \\
  owed' = \{ r : owed | rgen~r \neq curGen \} \\
  nreq' = nreq \\
  creq' = creq \\
  rgen' = rgen \\
  (cstate = cwait \implies cstate' = cidle) \\
  (cstate \neq cwait \implies cstate' = cstate) \\
  broken' = broken \cup \{ curGen \} \\
  deliberate' = deliberate \\
  crossGen' = crossGen \\
  crossReq' = crossReq \\
  lst' = lst \oplus \{ t? \mapsto lexit \} \\
  armed' = armed \oplus \{ t? \mapsto fclear \} \\
  lgen' = lgen \\
  bufr' = bufr \\
  (lst~t? = lstop \implies
    afterStop' = afterStop \oplus \{ t? \mapsto afterStop~t? + 1 \}) \\
  (lst~t? \neq lstop \implies afterStop' = afterStop)
\end{schema}

The reader's stream broke --- the read failed, the display closed it, or a frame did not decode --- and it reports by compare-and-close on its own generation.  If that generation is the open current one, it is closed and recorded as broken.

\subsection{BrokenStale}

\begin{schema}{BrokenStale}
  \Delta LL \\
  \Xi Conn \\
  \Xi Req \\
  \Xi Hist \\
  t? : T
\where
  armed~t? = fset \\
  \lnot (held = fset \land lgen~t? = curGen) \\
  lst' = lst \oplus \{ t? \mapsto lexit \} \\
  armed' = armed \oplus \{ t? \mapsto fclear \} \\
  lgen' = lgen \\
  bufr' = bufr \\
  (lst~t? = lstop \implies
    afterStop' = afterStop \oplus \{ t? \mapsto afterStop~t? + 1 \}) \\
  (lst~t? \neq lstop \implies afterStop' = afterStop)
\end{schema}

Otherwise the report is a no-op and the reader exits.  The current generation is never touched on behalf of another's reader.

\subsection{Halt}

\begin{schema}{Halt}
  \Xi LL
\where
  curGen = maxGen \\
  held = fclear \\
  cstate = cidle \\
  \forall u : T \spot lst~u = lexit
\end{schema}

The carrier is finite, so the run ends when every generation has been dialed and closed, every reader has exited and the caller is idle.  \texttt{Halt} is enabled in that quiescent state and nowhere else: a reachable state with no enabled operation that is \emph{not} this one is a deadlock.


\subsection{SendNoWait}

\begin{schema}{SendNoWait}
  \Delta LL \\
  \Xi Conn \\
  \Xi Thr \\
  \Xi Hist
\where
  cstate = cidle \\
  held = fset \\
  nreq < maxReq \\
  nreq' = nreq + 1 \\
  creq' = creq \\
  cstate' = cstate \\
  rgen' = rgen \oplus \{ nreq + 1 \mapsto curGen \} \\
  owed' = owed \cup \{ nreq + 1 \} \\
  qrep' = qrep
\end{schema}
A write that makes a reply owed and that nobody waits for: the replicator's \texttt{send\_scene}.  The display acks every scene (\texttt{render\_loop.\_handle\_scene}), so the ack arrives in the generation's sink with no waiter.

\end{document}
